% === Begin Label Data ===
% ID: 790
% 难度星级: 
% 标签: 
% 备注: 
% 组卷引用次数: 0
% === End  Label Data ===

\begin{problem}{2010}{G}{天津卷（文）}{15}{数列}
设 $\{a_n\}$ 是等比数列，公比 $q=\sqrt{2}$，$S_n$ 为 $\{a_n\}$ 的前 $n$ 项和．记 $T_n=\dfrac{17S_n-S_{2n}}{a_{n+1}}$，$n\in\mathbb{N}^*$．设 $T_{n_0}$ 为数列 $\{T_n\}$ 的最大项，则 $n_0=$\underline{\hspace{4em}}．
\end{problem}

\begin{answer}
$4$
\end{answer}

\begin{solutions}
利用等比数列的通项公式、前 $n$ 项和公式化简 $T_n=\dfrac{17S_n-S_{2n}}{a_{n+1}}$ 求解．

思路 1：由已知，得 $S_n=\dfrac{a_1(1-q^n)}{1-q}$，$S_{2n}=\dfrac{a_1(1-q^{2n})}{1-q}$，$a_{n+1}=a_1\cdot q^n$．

所以 $T_n=\dfrac{17S_n-S_{2n}}{a_{n+1}}=\dfrac{17(1-q^n)-(1-q^{2n})}{(1-q)\cdot q^n}=\dfrac{q^{2n}-17q^n+16}{(1-q)\cdot q^n}=\dfrac{1}{1-q}\left(q^n+\dfrac{16}{q^n}-17\right)$，

因为 $q^n+\dfrac{16}{q^n}\geqslant 8$，故 $T_n\leqslant\dfrac{-9}{1-\sqrt{2}}$ 当且仅当 $q^n=\dfrac{16}{q^n}$ 时等号成立．因为 $q=\sqrt{2}$，解得 $n=4$．即 $n_0=4$．

思路 2：$a_n=a_1(\sqrt{2})^{n-1}$，$S_n=\dfrac{a_1(1-(\sqrt{2})^n)}{1-\sqrt{2}}$，则
$$
\begin{aligned}
T_n&=\dfrac{17\dfrac{1-(\sqrt{2})^n}{1-\sqrt{2}}-\dfrac{1-(\sqrt{2})^{2n}}{1-\sqrt{2}}}{(\sqrt{2})^n}\\
&=\dfrac{1}{1-\sqrt{2}}\cdot\dfrac{16-17(\sqrt{2})^n+(\sqrt{2})^{2n}}{(\sqrt{2})^n}\\
&=\dfrac{1}{1-\sqrt{2}}\left(\dfrac{16}{(\sqrt{2})^n}+(\sqrt{2})^n-17\right).
\end{aligned}
$$
由基本不等式，$\dfrac{16}{(\sqrt{2})^n}+(\sqrt{2})^n\geqslant 2\sqrt{\dfrac{16}{(\sqrt{2})^n}\cdot(\sqrt{2})^n}=8$，等号成立条件是 $\dfrac{16}{(\sqrt{2})^n}=(\sqrt{2})^n$，即 $n=4$．因此，$T_n\leqslant T_4$．
\end{solutions}